\chapter{Ricci Flow with surgery: the definition}
\label{chap:ricci-flow-with-surgery}
\label{sect:surgery}



In this chapter we introduce Ricci flows with surgery. These objects
are closely related to generalized Ricci flows but they differ
slightly. The space-time of a Ricci flow with surgery has an open
dense subset that is a manifold, and the restriction of the Ricci
flow with surgery to this open subset is a generalized Ricci flow.
Still there are other, more singular points allowed in a Ricci flow
with surgery.

\section{Surgery space-time}


\begin{definition}[Space-time: paracompact Hausdorff space with a time function]
\label{def:space-time-general}
By a {\em space-time} we mean a paracompact Hausdorff space
${\mathcal M}$ with a continuous function ${\bf t}\colon {\mathcal
M}\to \Ar$, called {\em time}. We require that the image of ${\bf
t}$ be an interval $I$, finite or infinite with or without
endpoints, in $\Ar$. The interval $I$ is called the {\em time-interval of
definition of space-time}. The initial point of $I$, if there is
one,  is the {\em initial time} and the final point of $I$, if there
is one, is the {\em final time}. The level sets of ${\bf t}$ are
called the {\em time-slices} of space-time, and the preimage of the
initial (resp., final) point of $I$ is the {\em initial} (resp.,
{\em final}) time-slice.
\end{definition}

We are interested in a certain class of space-times, which we call
{\em surgery space-times}. These objects have a `smooth structure'
(even though they are not smooth manifolds). As in the case of a
smooth manifold, this smooth structure is given by local coordinate
charts with appropriate overlap functions.

\subsection{An exotic chart}

There is one exotic chart, and we begin with its description. To
define this chart we consider the open unit square $(-1,1)\times
(-1,1)$. We shall define a new topology, denoted by ${\mathcal P}$,
on this square. The open subsets of ${\mathcal P}$ are the open
subsets of the usual topology on the open square together with open
subsets of $(0,1)\times [0,1)$. Of course, with this topology the
`identity' map $\iota\colon {\mathcal P}\to (-1,1)\times (-1,1)$  is
a continuous map. Notice that the restriction of the topology of
${\mathcal P}$ to the complement of the closed subset $[0,1)\times
\{0\}$ is a homeomorphism onto the corresponding subset of the open
unit square. Notice that the complement of $(0,0)$ in ${\mathcal P}$
is a manifold with boundary, the boundary being $(0,1)\times\{0\}$.
(in the Introduction.)

Next, we define a `smooth structure' on ${\mathcal P}$ by defining a
sheaf of germs of `smooth' functions. The restriction of this sheaf
of germs of `smooth functions'  to the complement of $(0,1)\times
\{0\}$ in ${\mathcal P}$ is the usual sheaf of germs of smooth
functions on the corresponding subset of the open unit square.
In particular, a function is smooth near $(0,0)$ if and only if its restriction
to some neighborhood of $(0,0)$ is the pullback under $\iota$
of a usual smooth function on a
neighborhood of the origin in the square.
 Now
let us consider the situation near a point of the form $x=(a,0)$ for
some $0<a<1$. This point has arbitrarily small neighborhoods $V_n$
that are identified under $\iota$ with open subsets of $(0,1)\times
[0,1)$. We say that a function  $f$ defined in a neighborhood of $x$
in ${\mathcal P}$ is {\em smooth at $x$} if its restriction to one
of these neighborhoods $V_n$ is the pullback via $\iota|_{V_n}$ of a
smooth function in the usual sense on the open subset $\iota(V_n)$
of the upper half space. One checks directly that this defines a
sheaf of germs of `smooth' functions on ${\mathcal P}$. Notice that
the restriction of this sheaf to the complement of $(0,0)$ is the
structure sheaf of smooth functions of a smooth manifold with
boundary. Notice that the map $\iota\colon {\mathcal P}\to
(-1,1)\times (-1,1)$ is a smooth map in the sense that it pulls back
smooth functions on open subsets of the open unit square to smooth
functions on the corresponding open subset of ${\mathcal P}$.



Once we have the notion of smooth functions on ${\mathcal P}$, there
is the categorical notion of a diffeomorphism between open subsets
of ${\mathcal P}$: namely a homeomorphism with the property that it
and its inverse pull back smooth functions to smooth functions. Away
from the origin, this simply means that the map is a diffeomorphism
in the usual sense between manifolds with boundary, and in a
neighborhood of $(0,0)$ it factors through a diffeomorphism of
neighborhoods of the origin  in the square. While $\iota\colon
{\mathcal P}\to (-1,1)\times (-1,1)$ is a smooth map, it is not a
diffeomorphism.

We define the {\sl tangent bundle} of ${\mathcal P}$ in the usual
manner. The tangent space at a point is the vector space of
derivations of the germs of smooth functions at that point. Clearly,
away from $(0,0)$ this is the usual ($2$-plane) tangent bundle of
the smooth manifold with boundary. The germs of smooth functions at
$(0,0)$ are, by definition,  the pullbacks under $\iota$ of germs
of smooth functions at the origin for the unit square, so that the
tangent space of ${\mathcal P}$ at $(0,0)$ is identified  with the
tangent space of the open unit square at the origin. In fact, the
map $\iota$ induces an isomorphism from the tangent bundle of
${\mathcal P}$ to the pullback under $\iota$ of the tangent bundle
of the square. In particular, the tangent bundle of ${\mathcal P}$
has a given trivialization from the partial derivatives $\partial_x$
and $\partial_y$ in the coordinate directions on the square. We use
this trivialization to induce a smooth structure on the tangent
bundle of ${\mathcal P}$: that is to say, a section of $T{\mathcal
P}$ is smooth if and only if it can be written as $\alpha
\partial_x+\beta
\partial_y$ with $\alpha$ and $\beta $ being smooth functions on
${\mathcal P}$. The smooth structure agrees off of $(0,0)\in
{\mathcal P}$ with the usual smooth structure on the tangent bundle
of the smooth manifold with boundary.
By a {\em smooth vector field on ${\mathcal P}$}  we mean a smooth
section of the tangent bundle of ${\mathcal P}$.
 Smooth vector fields act as derivations on the smooth
functions on ${\mathcal P}$.

 We let ${\bf t}_{\mathcal P}\colon {\mathcal P}\to \Ar$ be the pullback via
$\iota$ of the usual projection to the second factor on the unit
square. We denote by $\chi_{\mathcal P}$ the smooth vector field
$\iota^*\partial_2$. Clearly, $\chi_{\mathcal P}({\bf t}_{\mathcal
P}) =1$. Smooth vector fields on ${\mathcal P}$ can be uniquely
integrated locally to smooth integral curves in ${\mathcal P}$. (At
a manifold with boundary point, of course only vector fields
pointing into the manifold can be locally integrated.)

\subsection{Coordinate charts for a surgery space-time}

Now we are ready to introduce the types of coordinate charts that we
shall use in our definition of a surgery space-time. Each coordinate
patch comes equipped with a smooth structure (a sheaf of germs of
smooth functions) and a tangent bundle with a smooth structure, so
that smooth vector fields act as derivations on the algebra of
smooth functions. There is also a distinguished smooth function,
denoted ${\bf t}$, and a smooth vector field, denoted $\chi$,
required to satisfy $\chi({\bf t})=1$. There are three types of
coordinates:
\begin{enumerate}
\item[(1)] The coordinate patch is an open subset of the strip ${\Ar}^{n}\times I$,
where $I$ is an interval, with its usual smooth structure and
tangent bundle; the function ${\bf t}$ is the projection onto $I$;
and the vector field $\chi$ is the unit tangent vector in the
positive direction tangent to the foliation with leaves $\{x\}\times
I$. The initial point of $I$, if there is one, is the initial time
of the space-time and the final point of $I$, if there is one, is
the final time of the space-time.
\item[(2)] The coordinate patch an open subset of $\Ar^n\times [a,\infty)$, for some $a\in \Ar$,
 with its
usual smooth structure as a manifold with boundary and its usual
smooth tangent bundle; the function ${\bf t}$ is the projection onto
the second factor; and the vector field is the coordinate partial
derivative associated with the second factor. In this case we
require that $a$ not be the initial time of the Ricci flow.
\item[(3)] The coordinate patch is a product of ${\mathcal P}$ with an open subset of
$\Ar^{n-1}$ with the smooth structure (i.e., smooth functions and
the smooth tangent bundle) being the product of the smooth structure
defined above on ${\mathcal P}$ with the usual smooth structure of
an open subset of $\Ar^{n-1}$; the function ${\bf t}$ is, up to an
additive constant, the pullback of the function ${\bf t}_{\mathcal
P}$ given above on ${\mathcal P}$; and the vector field $\chi$ is
the image of the vector field $\chi_{\mathcal P}$ on ${\mathcal P}$,
given above, under the product decomposition.
\end{enumerate}

An ordinary Ricci flow is covered by coordinate charts of the first
type.  The second and third are two extra types of coordinate charts
for a Ricci flow with surgery that are not allowed in generalized
Ricci flows. Charts of the second kind are smooth
manifold-with-boundary charts, where the boundary is contained in a
single time-slice,  not the initial time-slice, and the flow exists
for some positive amount of forward time from this manifold.


All the structure described above for ${\mathcal P}$ -- the smooth
structure, the tangent bundle with its smooth structure, smooth
vector fields acting as derivations on smooth functions -- exist for
charts of the third type. In addition, the unique local integrability
of smooth vector fields  hold for coordinate charts of the third
type. Analogous results for coordinate charts of the first two types
are clear.

Now let us describe the allowable overlap functions between charts.
Between charts of the first and second type these are the smooth
overlap functions in the usual sense that preserve the functions
${\bf t}$ and the vector fields $\chi$ on the patches. Notice that
because the boundary points in charts of the second type are
required to be at times other than the initial and final times, the
overlap of a chart of type one and a chart of type two is disjoint
from the boundary points of each. Charts of the first two types are
allowed to meet a chart of the third type only in its manifold and
manifold-with-boundary points. For overlaps between charts of the
first two types with a chart of the third type, the overlap
functions are  diffeomorphisms between open subsets preserving the
local time functions ${\bf t}$ and the local vector fields $\chi$.
Thus, all overlap functions are diffeomorphisms in the sense given
above.

\subsection{Definition and basic properties of surgery space-time}

\begin{definition}[Surgery space-time: smooth, exposed, and singular points]
\label{def:surgery-space-time}
\uses{def:space-time-general}\label{spacetime}
A {\em surgery space-time} is a
space-time  ${\mathcal M}$ equipped with a maximal atlas  of charts
covering ${\mathcal M}$, each chart being of one of the three types
listed above, with the overlap functions being diffeomorphisms
preserving the functions ${\bf t}$ and the vector fields $\chi$. The
points with neighborhoods of the first type are called {\em smooth
points}, those with neighborhoods of the second type but not the
first type are called {\em exposed points},
and all the other points are called {\em singular
points}. Notice that the union of
the set of smooth points and the set of exposed points forms a
smooth manifold with boundary (possibly disconnected). Each
component of the boundary of this manifold is contained in a single
time-slice. The union of those components contained in a time
distinct from the initial time and the final time is called the {\em
exposed region}. and the boundary points of
the closure of the exposed region form the set of the singular
points of ${\mathcal M}$. (Technically, the exposed points are
singular, but we reserve this word for the most singular points.) An
$(n+1)$-dimensional surgery space-time is by definition of
homogeneous dimension $n+1$.

By construction, the local smooth functions ${\bf t}$ are compatible
on the overlaps and hence fit together to define a global smooth
function ${\bf t}\colon {\mathcal M}\to \Ar$, called the {\em time}
function. The level sets of this function are called the {\em
time-slices} of the space-time, and ${\bf t}^{-1}(t)$ is denoted
$M_t$. Similarly, the tangent bundles of the various charts are
compatible under the overlap diffeomorphisms and hence glue together
to give a global smooth tangent bundle on space-time.  The smooth
sections of this vector bundle, the smooth vector fields on space
time, act as derivations on the smooth functions on space-time. The
tangent bundle of an $(n+1)$-dimensional surgery space-time is a
vector bundle of dimension $(n+1)$. Also, by construction the local
vector fields $\chi$ are compatible and hence glue together to
define a global vector field, denoted $\chi$. The vector field and
time function satisfy
$$\chi({\bf t})=1.$$
At the manifold points (including the exposed points) it is a usual
vector field. Along the exposed region and the initial time-slice
the vector field points into the manifold; along the final
time-slice it points out of the manifold.
\end{definition}




\begin{definition}[Embedding compatible with time and the vector field; maximal flow line]
\label{def:compatible-with-time-vector-field}
\uses{def:surgery-space-time}
Let ${\mathcal M}$ be a surgery space-time. Given a space $K$ and an
interval $J\subset \Ar$  we say that an embedding $K\times J\to
{\mathcal M}$ is {\em compatible with  time  and the vector field}
if: (i) the restriction of ${\bf t}$ to the image agrees with the
projection onto the second factor and (ii) for each $x\in X$ the
image of $\{x\}\times J$ is the integral curve for the vector field
$\chi$.
 If in addition $K$ is a subset of $M_{t}$
we require that $t\in J$ and that the map $K\times\{t\}\to M_{t}$ be
the identity. Clearly, by the uniqueness of integral curves for vector fields, two such
embeddings agree on their common interval of definition, so that,
given $K\subset M_{t}$ there is a maximal interval $J_K$ containing
$t$ such that such an embedding is defined on $K\times J_K$. In the
special case when $K=\{x\}$ for a point $x\in M_{t}$ we say that
such an embedding is {\em the maximal flow line} through $x$. The
embedding of the maximal interval through $x$ compatible with time
and the vector field $\chi$ is called {\em the domain of definition}
of the flow line through $x$. For a more general subset $K\subset
M_{t}$ there is an embedding $K\times J$ compatible with time and
the vector field $\chi$ if and only if, for every $x\in K$, the
interval $J$ is contained in the domain of definition of the flow
line through $x$.
\end{definition}


\begin{definition}[Regular and singular times of a surgery space-time]
\label{def:regular-singular-time-surgery}
\uses{def:surgery-space-time,def:compatible-with-time-vector-field}
Let ${\mathcal M}$ be a surgery space-time with $I$ as its time
interval of definition. We say that $t\in I$ is a {\em regular} time
if there is an interval $J\subset I$ which is an open neighborhood
in $I$ of $t$, and a diffeomorphism $M_{t}\times J\to {\bf
t}^{-1}(J)\subset{\mathcal M}$ compatible with time and the vector
field. A time is {\em singular} if it is not regular. Notice that if
all times are regular, then space-time is a product $M_t\times I$
with ${\bf t}$ and $\chi$ coming from the second factor.
\end{definition}



\begin{lemma}[Time-slices are smooth manifolds with a complementary horizontal subbundle]
\label{lem:time-slice-smooth-manifold}
\uses{def:surgery-space-time}
Let ${\mathcal M}$ be an $(n+1)$-dimensional surgery space-time, and fix $t$.
The restriction of the smooth structure on ${\mathcal M}$ to the time-slice
$M_{t}$ induces the structure of a smooth $n$-manifold on this time-slice. That
is to say, we have a smooth embedding of $M_{t}\to {\mathcal M}$. This smooth
embedding identifies the tangent bundle of $M_{t}$ with a codimension-one
subbundle of the restriction  of tangent bundle of ${\mathcal M}$ to $M_{t}$.
This subbundle is complementary to the line field spanned by $\chi$. These
codimension-one subbundles along the various time-slices fit together to form
a smooth, codimension-one subbundle of the tangent bundle of space-time.
 \end{lemma}


\begin{proof}
These statements are immediate for any coordinate patch, and hence
are true globally.
\end{proof}

\begin{definition}[Horizontal subbundle of a surgery space-time]
\label{def:horizontal-subbundle}
\uses{lem:time-slice-smooth-manifold}
We call the codimension-one subbundle of the tangent bundle of
${\mathcal M}$ described in the previous lemma  the {\em horizontal
subbundle}, and we denote it
${\mathcal HT}({\mathcal M})$.
\end{definition}

\section{The generalized Ricci flow equation}

In this section we introduce the Ricci flow equation for surgery
space-times, resulting in an object that we call Ricci flow with
surgery.

\subsection{Horizontal metrics}

\begin{definition}[Horizontal metric and its curvatures on a surgery space-time]
\label{def:horizontal-metric}
\uses{def:horizontal-subbundle}
By a {\em horizontal metric $G$ on a surgery space-time ${\mathcal
M}$} we mean a $C^\infty$ metric on ${\mathcal HT}{\mathcal M}$. For
each $t$, the horizontal metric $G$ induces a Riemannian metric,
denoted $G(t)$, on the time-slice $M_t$. Associated to a horizontal
metric $G$ we have the {\em horizontal covariant derivative},
denoted $\nabla$. This is a pairing between horizontal vector fields
$$X\otimes Y\mapsto \nabla_XY.$$
On each time slice $M_t$ it is the usual Levi-Civita connection associated to
the Riemannian metric $G(t)$. Given a function $F$ on space-time, by its
gradient $\nabla F$ we mean its horizontal gradient. The value of this gradient
at a point $q\in M_t$ is the usual $G(t)$-gradient of $F|_{M_t}$. In
particular, $\nabla F$ is a smooth horizontal vector field on space-time. The
horizontal metric $G$ on space-time has its (horizontal) curvatures $\Rm_G$.
These are smooth symmetric endomorphisms of the second exterior power of
${\mathcal HT}{\mathcal M}$. The value of $\Rm_G$ at a point $q\in M_t$ is
simply the usual Riemann curvature operator of $G(t)$ at the point $q$.
Similarly, we have the (horizontal) Ricci curvature $\Ric=\Ric_G$, a
section of the symmetric square of the horizontal cotangent bundle, and the
(horizontal) scalar curvature denoted $R=R_G$.
\end{definition}

The only reason for working in ${\mathcal HT}{\mathcal M}$ rather
than individually in each slice is to emphasize the fact that all
these horizontal quantities vary smoothly over the surgery
space-time.


Suppose that $t\in I$ is not the final time and suppose that
$U\subset M_{t}$ is an open subset with compact closure. Then there
is $\epsilon>0$ and an embedding $i_U\colon
U\times[t,t+\epsilon)\subset {\mathcal M}$ compatible with time and
the vector field. Of course, two such embeddings agree on their
common domain of definition. Notice also that for each $t'\in
[t,t+\epsilon)$ the restriction of the map $i_U$ to $U\times \{t'\}$
induces an diffeomorphism from $U$ to an open subset $U_{t'}$ of
$M_{t'}$. It follows that the local flow generated by the vector
field $\chi$ preserves the horizontal subbundle. Hence, the vector
field $\chi$ acts by Lie derivative on the sections of ${\mathcal
HT}({\mathcal M})$ and on all associated bundles (for example the
symmetric square of the dual bundle).




\subsection{The equation}



\begin{definition}[Ricci flow with surgery: the generalized Ricci flow equation]
\label{def:ricci-flow-with-surgery}
\uses{def:horizontal-metric,def:surgery-space-time}
A {\em  Ricci flow with surgery} is a pair $({\mathcal M},G)$
consisting of a surgery space-time ${\mathcal M}$ and a horizontal
metric $G$ on ${\mathcal M}$ such that for every $x\in {\mathcal M}$
 we have
\begin{equation}\label{GRFEq}
{\mathcal L}_\chi(G)(x)=-2\Ric_G(x)) \end{equation} as sections
of the symmetric square of the dual to ${\mathcal HT}({\mathcal
M})$. If space-time is $(n+1)$-dimensional, then we say that the
Ricci flow with surgery is $n$-dimensional (meaning of course that
each time-slice is an $n$-dimensional manifold).
\end{definition}

\begin{remark}[Lie derivative is one-sided at exposed points and initial/final times]
\label{rem:one-sided-lie-derivative}
\uses{def:ricci-flow-with-surgery}
 Notice that at an exposed point and at points at the initial and the final time
   the Lie derivative is a one-sided derivative.
\end{remark}


\subsection{Examples of Ricci flows with surgery}

\begin{example}[An ordinary Ricci flow is a Ricci flow with surgery]
\label{ex:ordinary-flow-as-surgery-flow}
\uses{def:ricci-flow-with-surgery}
One example of a Ricci flow with surgery is ${\mathcal M}=M_0\times
[0,T)$ with time function ${\bf t}$ and the vector field $\chi$
coming from the second factor. In this case the Lie derivative
${\mathcal L}_\chi$ agrees with the usual partial derivative in the
time direction, and hence our generalized Ricci flow equation is the
usual Ricci flow equation. This shows that an ordinary Ricci flow is indeed
a Ricci flow with surgery.
\end{example}

The next lemma gives an example of a  Ricci flow with surgery where
the topology of the time-slices changes.


\begin{lemma}[Gluing two Ricci flows along a common time-slice yields a Ricci flow with surgery]
\label{lem:gluing-ricci-flows-with-surgery}
\uses{def:ricci-flow-with-surgery}\label{gluingsoln}
Suppose that we have manifolds $M_1\times (a,b]$ and $M_2\times
[b,c)$ and compact, smooth codimension-$0$ submanifolds
$\Omega_1\subset M_1$ and $\Omega_2\subset M_2$ with open
neighborhoods $U_1\subset M_1$ and $U_2\subset M_2$ respectively.
Suppose we have a diffeomorphism $\psi\colon U_1\to U_2$ carrying
$\Omega_1$ onto $\Omega_2$. Let $(M_1\times (a,b])_0$ be the subset
obtained by removing $(M_1\setminus \Omega_1)\times \{b\}$ from
$M_1\times (a,b]$. Form the topological space
$${\mathcal M}=(M_1\times (a,b])_0\cup M_2\times [b,c)$$
where  $\Omega_1\times\{b\}$ in $(M_1\times (a,b])_0$ is identified
with $\Omega_2\times \{b\}$ using the restriction of $\psi$ to
$\Omega_1$. Then ${\mathcal M}$ naturally inherits the structure of
a surgery space-time where the time function restricts to
$(M_1\times (a,b])_0$ and to $M_2\times [b,c)$ to be the projection
onto the second factor and the vector field $\chi$ agrees with the
vector fields coming from the second factor on each of $(M_1\times
(a,b])_0$ and $M_2\times [b,c)$.

 Lastly, given Ricci
flows $(M_1,g_1(t)),\ a<t\le b$, and $(M_2,g_2(t)),\ b\le t<c$, if
$\psi\colon (U_1,g_1(b))\to (U_2,g_2(b))$ is an isometry, then these
families fit together to form a smooth horizontal metric $G$ on
${\mathcal M}$ satisfying the Ricci flow equation, so that
$({\mathcal M},G)$ is a  Ricci flow with surgery.
\end{lemma}

\begin{proof}
\uses{prop:smooth-patching}
As the union of Hausdorff spaces along closed subsets, ${\mathcal
M}$ is a Hausdorff topological space. The time function is the one
induced from the projections onto the second factor. For any point
outside the $b$ time-slice there is the usual smooth coordinate
coming from the smooth manifold $M_1\times (a,b)$ (if $t<b$) or
$M_2\times (b,c)$ (if $t>b$). At any point of
$(M_2\setminus\Omega_2)\times \{b\}$ there is the smooth manifold
with boundary coordinate patch coming from $M_2\times[b,c)$. For any
point in $\mathrm{int}(\Omega_1)\times\{b\}$ we have the smooth
manifold structure obtained from gluing $(\mathrm{int}(\Omega_1))\times
(a,b]$ to $\mathrm{int}(\Omega_2)\times [b,c)$ along the $b$ time-slice
by $\psi$. Thus, at all these points we have neighborhoods on which
our data determine a smooth manifold structure. Lastly, let us
consider a point $x\in
\partial \Omega_1\times \{b\}$.
Choose local coordinates $(x^1,\ldots,x^n)$ for a neighborhood $V_1$
of $x$ such that $\Omega_1\cap V_1=\{x^n\le 0\}$. We can assume that
$\psi$ is defined on all of $V_1$. Let $V_2=\psi(V_1)$ and take the
 local coordinates on $V_2$ induced from the $x^i$ on $V_1$.
 Were we to identify $V_1\times (a,b]$
with $V_2\times [b,c)$ along the $b$ time-slice using this map, then
this union would be a smooth manifold.  There is a
neighborhood of the point $(x,b)\in {\mathcal M}$ which is obtained
from the smooth manifold $V_1\times (a,b]\cup_\psi V_2\times [b,c)$
by inducing a new topology where the open subsets are, in addition
to the usual ones, any open subset of the form $\{x^n>0\}\times
[b,b')$ where $b<b'\le c$. This then gives the coordinate charts of
the third type near the points of $\partial \Omega_2\times \{b\}$.
Clearly, since the function ${\bf t}$ and the vector field
$\partial/\partial t$ are smooth on $V_1\times (a,b]\cup_\psi
V_2\times [b,c)$, we see that these objects glue together to form
smooth objects on ${\mathcal M}$.



Given the Ricci flows $g_1(t)$ and $g_2(t)$ as in the statement,
they clearly determine a (possibly singular) horizontal metric on
${\mathcal M}$. This horizontal metric is clearly smooth except
possibly along the $b$ time-slice. At any point of
$(M_2\setminus\Omega_2)\times\{b\}$ we have a one-sided smooth
family, which means that on this set the horizontal metric is
smooth. At a point of $\mathrm{int}(\Omega_2)\times \{b\}$, the fact
that the metrics fit together smoothly is an immediate consequence
of Proposition~\ref{patch}. At a point $x\in
\partial \Omega_2\times \{b\}$ we have neighborhoods $V_2\subset M_2$
of $x$ and $V_1\subset M_1$ of $\psi^{-1}(x)$ that are isometrically
identified by $\psi$. Hence, again by Lemma~\ref{patch} we see that
the Ricci flows fit together to form a smooth family of metrics on
$V_1\times (a,b]\cup_\psi V_2\times [b,c)$. Hence, the induced
horizontal metric on ${\mathcal M}$ is smooth near this point.
\end{proof}

The following is obvious from the definitions.

\begin{proposition}[Smooth and exposed points of a Ricci flow with surgery form a generalized Ricci flow]
\label{prop:interior-generalized-ricci-flow}
\uses{def:ricci-flow-with-surgery,def:generalized-ricci-flow}
Suppose that $({\mathcal M},G)$ is a Ricci flow with surgery. Let
$\mathrm{int}{\mathcal M}$ be the open subset consisting of all
smooth $(n+1)$-manifold points, plus all manifold-with-boundary
points at the initial time and the final time. This space-time
inherits the structure of a smooth manifold with boundary. This
structure together with the restrictions to it of ${\bf t}$ and the
vector field $\chi$ and the restriction of the horizontal metric $G$
form a generalized Ricci flow whose underlying smooth manifold is $\mathrm{int}{\mathcal M}$.
\end{proposition}


\subsection{Scaling and translating Ricci flows with surgery}

Suppose that $({\mathcal M},G)$ is a Ricci flow with surgery. Let
$Q$ be a positive constant. Then we can define a new  Ricci flow
with surgery by setting $G'=QG$, ${\bf t'}=Q{\bf t}$ and
$\chi'=Q^{-1}\chi$. It is easy to see that the resulting data still
satisfies the generalized Ricci flow equation,
Equation~(\ref{GRFEq}). We denote this new  Ricci flow with surgery by
$(Q{\mathcal M},QG)$ where the changes in ${\bf t}$ and $\chi$ are
indicated by the factor $Q$ in front of the space-time.

It is also possible to translate a Ricci flow with surgery $({\mathcal
M},G)$ by replacing the time function ${\bf t}$ by ${\bf t'}={\bf
t}+a$ for any constant $a$, and leaving $\chi$ and $G$ unchanged.


\subsection{More basic definitions}



\begin{definition}[Metric ball in a time-slice of space-time]
\label{def:metric-ball-spacetime}
\uses{def:surgery-space-time}
Let $({\mathcal M},G)$ be a Ricci flow with surgery, and let $x$ be
a point of space-time. Set $t={\bf t}(x)$. For any  $r>0$ we define
$B(x,t,r)\subset M_t$ to be the metric ball of radius $r$ centered
at $x$ in the Riemannian manifold $(M_t,G(t))$.
\end{definition}

\begin{definition}[Forward and backward parabolic neighborhoods in a Ricci flow with surgery]
\label{def:parabolic-neighborhood-surgery}
\uses{def:metric-ball-spacetime,def:compatible-with-time-vector-field}
Let $({\mathcal M},G)$ be a Ricci flow with surgery, and let $x$ be
a point of space-time. Set $t={\bf t}(x)$. For any $r>0$ and $\Delta
t>0$ we say that the {\em backward parabolic neighborhood}
$P(x,t,r,-\Delta t)$ {\em exists} in ${\mathcal M}$ if there is an
embedding $B(x,t,r)\times (t-\Delta t,t]\to {\mathcal M}$ compatible
with time and the vector field. Similarly, we say
that the {\em forward parabolic neighborhood} $P(x,t,r,\Delta t)$
{\em exists} in ${\mathcal M}$ if there is an embedding
$B(x,t,r)\times [t,t+\Delta t)\to {\mathcal M}$ compatible with time
and the vector field. A {\em parabolic neighborhood} is either a
forward or backward parabolic neighborhood.
\end{definition}


\begin{definition}[Kappa-non-collapsed on scales at most r_0, for surgery flows]
\label{def:kappa-noncollapsed-surgery}
\uses{def:parabolic-neighborhood-surgery}\label{defnkappa}
Fix $\kappa>0$ and $r_0>0$. We say that a Ricci flow with surgery
$({\mathcal M},G)$ is {\em $\kappa$-noncollapsed on scales $\le
r_0$} if the following holds for every point $x\in {\mathcal M}$ and
for every $r\le r_0$. Denote ${\bf t}(x)$ by $t$. If the  parabolic
neighborhood $P(x,t,r,-r^2)$ exists in ${\mathcal M}$ and if
$|\Rm_G|\le r^{-2}$ on $P(x,t,r,-r^2)$, then  $\Vol\,B(x,t,r)\ge \kappa r^3$.
\end{definition}


\begin{remark}[Non-collapsing constant for epsilon-round components depends on fundamental group order]
\label{rem:epsilon-round-kappa-dependence}
\uses{def:kappa-noncollapsed-surgery,def:epsilon-round-component}
For $\epsilon>0$ sufficiently small,
an $\epsilon$-round component satisfies the first condition in the above definition
 for some $\kappa>0$ depending only on
the order of the fundamental group of the underlying manifold, but there is no universal
$\kappa>0$ that works for all $\epsilon$-round manifolds. Fixing an integer $N$ let ${\mathcal C}_N$
be the class of closed $3$-manifolds
 with the property that any finite free factor of $\pi_1(M)$
has order at most $N$. Then any $\epsilon$-round component of any
time-slice of any Ricci flow $({\mathcal M},G)$ whose initial
conditions consist of a manifold in ${\mathcal C}_N$ will have
fundamental group of order at most $N$ and hence will satisfy the first condition
in the above definition for some $\kappa>0$ depending only on $N$.
\end{remark}


We also have the notion of the curvature being  pinched toward
positive, analogous to the notions for Ricci flows and  generalized Ricci flows.

\begin{definition}[Curvature pinched toward positive, for a Ricci flow with surgery]
\label{def:curvature-pinched-toward-positive-surgery}
\uses{def:ricci-flow-with-surgery,def:curvature-operator}
Let $({\mathcal M},G)$ be a $3$-dimensional Ricci flow with surgery,
whose time domain of definition is contained in $[0,\infty)$. For
any $x\in {\mathcal M}$ we denote the eigenvalues of $\Rm(x)$ by
$\lambda(x)\ge \mu(x)\ge \nu(x)$ and we set $X(x)=\mathrm{max}(0,-\nu(x))$.
 We say
that its {\em curvature is pinched toward positive} if the following
hold for every $x\in {\mathcal M}$:
\begin{enumerate}
\item[(1)] $R(x)\ge \frac{-6}{1+4{\bf t}(x)}$.
\item[(2)] $R(x)\geq 2X(x)\left(\mathrm{log}X(x)+\mathrm{log}(1+{\bf t}(x))-3\right)$, whenever $0<X(x)$.
\end{enumerate}

Let $(M,g)$ be a Riemannian manifold and let $T\ge 0$.
We say that $(M,g)$ has curvature {\em pinched toward positive up to time $T$} if
the above two inequalities hold for all $x\in M$ with ${\bf t}(x)$ replaced by $T$.
\end{definition}

Lastly, there is the definition of canonical neighborhoods for a
Ricci flow with surgery, there is the following extension of the
notion for a generalized Ricci flow.


\begin{definition}[Strong (C,epsilon)-canonical neighborhood assumption with parameter r]
\label{def:strong-canonical-neighborhood-assumption}
\uses{def:evolving-epsilon-neck,def:canonical-neighborhood,def:strong-epsilon-neck-canonical-neighborhood,def:compatible-with-time-vector-field}
Fix constants $(C,\epsilon)$ and a constant $r$. We say that a Ricci
flow with surgery $({\mathcal M},G)$ {\em satisfies the strong
$(C,\epsilon)$-canonical neighborhood assumption with parameter $r$}
if every point $x\in {\mathcal M}$ with $R(x)\ge r^{-2}$ has a strong
$(C,\epsilon)$-canonical neighborhood in ${\mathcal M}$. In all
cases except that of the strong $\epsilon$-neck, the strong
canonical neighborhood of $x$ is a subset of the time-slice
containing $x$, and the notion of a $(C,\epsilon)$-canonical
neighborhood has exactly the same meaning as in the case of an
ordinary Ricci flow. In the case of a strong $\epsilon$-neck
centered at $x$ this means that there is an embedding
$\left(S^2\times (-\epsilon^{-1},\epsilon^{-1})\right)\times ({\bf
t}(x)-R(x)^{-1},{\bf t}(x)]\to {\mathcal M}$, mapping $(q_0,0)$ to $x$,
 where $q_0$ is the basepoint of $S^2$, an embedding
compatible with time and the vector field,
such that the pullback
of $G$ is a Ricci flow on $S^2\times (-\epsilon^{-1},\epsilon^{-1})$
which, when the time is shifted by $-{\bf t}(x)$ and then the flow
is rescaled by $R(x)$, is within $\epsilon$ in the
$C^{[1/\epsilon]}$-topology of the standard evolving round cylinder
$\left(S^2\times(-\epsilon^{-1},\epsilon^{-1}),h_0(t)\times
ds^2\right),\ -1<t\le 0$, where the scalar curvature of the $h_0(t)$
is $1-t$.
\end{definition}

Notice that $x$ is an exposed point or sufficiently close to an exposed point
then $x$ cannot be the center of a strong $\epsilon$-neck.

